% ============================================================================
% Chapter 7: The Software and the v4.0 Roadmap
% Study Guide — Cardano Credential Verification
% ============================================================================

\chapter{The Software and the v4.0 Roadmap}
\label{ch:software-roadmap}

This chapter is divided into three parts.  Part~A describes the
\texttt{cardano-license} Python package---its architecture, API surface,
database schema, and testing infrastructure.  Part~B presents the v4.0
roadmap: Aiken smart contracts, CIP-33 reference scripts, validity-token
consolidation, HSM key storage, and KERI/multi-sig governance.  Part~C
collects a comprehensive glossary of every technical term used throughout the
study guide.

% ── TikZ style definitions for this chapter ──────────────────────────────────
\tikzset{
  guidebox/.style={draw, rounded corners=4pt, thick, align=center,
    font=\sffamily\small, fill=white, minimum height=1cm},
  guidearrow/.style={-{Latex[length=2.5mm]}, thick, color=gray!70!black},
  guidelabel/.style={font=\sffamily\footnotesize, text=gray!50!black},
}


% ============================================================================
%  PART A — THE SOFTWARE
% ============================================================================

\vspace{0.5em}
\noindent\rule{\textwidth}{0.8pt}
\begin{center}
  {\Large\sffamily\bfseries Part A --- The Software}
\end{center}
\noindent\rule{\textwidth}{0.4pt}
\vspace{0.5em}


% ----------------------------------------------------------------------------
\section{Package Overview}
\label{sec:pkg-overview}

The \texttt{cardano-license} package is an open-source Python library
released under the MIT license.  As of this writing, the current release is
version~2.1.0.  The package provides a complete toolkit for professional
credential management on the Cardano blockchain, encompassing wallet
generation, license NFT minting, signature-token operations, validity
management, CIP-68 reference-token revocation, Plutus~V2 smart-contract
deployment, and dues enforcement.

Installation is straightforward via the Python Package Index:

\begin{lstlisting}[language=bash, caption={Installing cardano-license}]
pip install cardano-license
\end{lstlisting}

The package requires Python~3.10 or later and depends on four core
libraries.  \textbf{PyCardano} serves as the Cardano SDK, providing
HD~wallet generation, transaction building, Plutus~V2 script support, and
address derivation.  \textbf{Blockfrost} provides the chain-query API,
enabling the package to interact with the Cardano blockchain without running
a full node.  \textbf{aiosqlite} provides asynchronous SQLite database
access, allowing the package to maintain local state (wallets, licenses,
signatures, contracts) without blocking the event loop.
\textbf{cryptography} provides AES-256-GCM authenticated encryption for
wallet key files, using PBKDF2-HMAC-SHA256 key derivation with 100\,000
iterations.

\begin{table}[htbp]
\centering
\caption{Core dependencies of the \texttt{cardano-license} package.}
\label{tab:dependencies}
\begin{tabular}{@{}llp{6.5cm}@{}}
\toprule
\textbf{Library} & \textbf{Role} & \textbf{Usage} \\
\midrule
PyCardano        & Cardano SDK   & HD wallets, transaction building, Plutus~V2 scripts, CIP-1852 derivation \\
Blockfrost       & Chain API     & UTxO queries, transaction submission, balance lookups \\
aiosqlite        & Async DB      & Local persistence of wallets, licenses, signatures, contracts \\
cryptography     & Encryption    & AES-256-GCM key-file encryption with PBKDF2 derivation \\
\bottomrule
\end{tabular}
\end{table}

The package follows a layered design philosophy: all on-chain interactions
flow through PyCardano and Blockfrost, all local state is managed through
aiosqlite, and all key material is protected by the cryptography module.
This separation of concerns means that swapping the chain backend (e.g.,
from Blockfrost to Ogmios) or the database engine (e.g., from SQLite to
PostgreSQL) would require changes to only one layer.


% ----------------------------------------------------------------------------
\section{Architecture}
\label{sec:architecture}

The package is organized into four Python modules, each with a clearly
defined responsibility.  Figure~\ref{fig:module-diagram} illustrates the
dependency relationships between modules and external libraries.

\begin{figure}[htbp]
\centering
\begin{tikzpicture}[node distance=1.4cm and 2.2cm]
  % ── Internal modules ──
  \node[guidebox, fill=blue!8, minimum width=4.2cm]
    (core) {\textbf{core.py}\\{\scriptsize\textsf{Wallets, Minting, Signing}}\\{\scriptsize\textsf{Validity, CIP-68, Dues}}\\{\scriptsize\textsf{$\sim$4\,700 lines}}};

  \node[guidebox, fill=green!8, minimum width=3.8cm, right=2.4cm of core]
    (tx) {\textbf{tx\_utils.py}\\{\scriptsize\textsf{Fee Estimation, Tx Build}}\\{\scriptsize\textsf{UTxO Selection, Submit}}\\{\scriptsize\textsf{$\sim$800 lines}}};

  \node[guidebox, fill=orange!8, minimum width=3.8cm, below=1.6cm of core]
    (schema) {\textbf{schema.py}\\{\scriptsize\textsf{10 Tables, init\_db()}}\\{\scriptsize\textsf{DDL Statements}}};

  \node[guidebox, fill=purple!8, minimum width=3.8cm, below=1.6cm of tx]
    (init) {\textbf{\_\_init\_\_.py}\\{\scriptsize\textsf{Public API Exports}}\\{\scriptsize\textsf{Version 2.1.0}}};

  % ── External libraries ──
  \node[guidebox, fill=yellow!8, minimum width=2.4cm, above=1.6cm of core]
    (pycardano) {\textbf{PyCardano}\\{\scriptsize\textsf{Cardano SDK}}};

  \node[guidebox, fill=yellow!8, minimum width=2.4cm, above=1.6cm of tx]
    (blockfrost) {\textbf{Blockfrost}\\{\scriptsize\textsf{Chain API}}};

  \node[guidebox, fill=red!6, minimum width=2.4cm, left=1.2cm of schema]
    (crypto) {\textbf{crypto.py}\\{\scriptsize\textsf{AES-256-GCM}}};

  \node[guidebox, fill=gray!10, minimum width=2.4cm, right=1.2cm of init]
    (aiosqlite) {\textbf{aiosqlite}\\{\scriptsize\textsf{Async SQLite}}};

  % ── Arrows ──
  \draw[guidearrow] (core) -- (pycardano);
  \draw[guidearrow] (core) -- (blockfrost);
  \draw[guidearrow] (tx)   -- (pycardano);
  \draw[guidearrow] (tx)   -- (core);
  \draw[guidearrow] (core) -- (schema);
  \draw[guidearrow] (core) -- (crypto);
  \draw[guidearrow] (schema) -- (aiosqlite);
  \draw[guidearrow] (init) -- (core);
  \draw[guidearrow] (init) -- (tx);
  \draw[guidearrow] (init) -- (schema);
\end{tikzpicture}
\caption{Module dependency diagram for \texttt{cardano-license}.  Arrows
  point from the dependent module to the dependency.  Yellow boxes are
  external libraries; colored boxes are internal modules.}
\label{fig:module-diagram}
\end{figure}

\subsection{core.py --- The Main Module}

The \texttt{core.py} module is the heart of the package, weighing in at
approximately 4\,700~lines of Python.  It implements every domain operation
in the credential lifecycle: HD~wallet generation following the CIP-1852
derivation standard, authority-restricted minting-policy creation,
CIP-25--compliant license NFT minting, CIP-68 reference-token management
for authority-controlled revocation, signature-token minting and transfer,
validity-token minting and renewal, document signing with on-chain metadata,
multi-signer work-product coordination, Plutus~V2 signature-collection
validators with atomic finalization, and dues-enforcement contracts with
grace-period logic.

The module is organized into logical sections, each corresponding to a phase
of the credential lifecycle.  Wallet operations come first, followed by
chain-context helpers, minting-policy management, license NFT operations,
signature-token operations, validity-token operations, CIP-68 revocation,
document signing and verification, work-product management, Plutus~V2
validators, and dues enforcement.  All database-touching functions are
\texttt{async} and use \texttt{aiosqlite} for non-blocking I/O.

\subsection{tx\_utils.py --- Transaction Utilities}

The \texttt{tx\_utils.py} module contains approximately 800~lines of
transaction-building utilities that abstract away the lower-level details of
Cardano transaction construction.  Its primary exports include fee
estimation (\texttt{estimate\_fee}, \texttt{estimate\_fee\_from\_context}),
UTxO selection (\texttt{select\_utxos}), and high-level transaction builders
(\texttt{build\_mint\_tx}, \texttt{build\_transfer\_tx},
\texttt{build\_multisig\_tx}).  It also provides
\texttt{submit\_tx} for transaction submission and
\texttt{wait\_for\_confirmation} for polling until a transaction appears
on-chain.

The module defines two important data classes.  \texttt{TxResult} captures
the outcome of any transaction operation---hash, fee, input/output counts,
minted assets, confirmation status, and any error.  \texttt{UTxOSelection}
encapsulates the result of UTxO selection, including the selected inputs,
total value, and change amount.

\subsection{schema.py --- Database Schema}

The \texttt{schema.py} module defines the SQLite schema for the package's
local state.  It contains ten \texttt{CREATE TABLE IF NOT EXISTS} statements
and exposes a single public function, \texttt{init\_db()}, which ensures all
tables exist.  The ten tables are described in detail in
Section~\ref{sec:db-schema}.

\subsection{\_\_init\_\_.py --- Package Exports}

The \texttt{\_\_init\_\_.py} module re-exports the public API from
\texttt{core.py}, \texttt{tx\_utils.py}, and \texttt{schema.py}.  It also
defines the package version (\texttt{\_\_version\_\_ = "2.1.0"}) and the
\texttt{\_\_all\_\_} list, which enumerates every public symbol.  This
module serves as the single entry point for consumers:
\texttt{from cardano\_license import generate\_wallet, mint\_license\_nft}.


% ----------------------------------------------------------------------------
\section{Key API Functions}
\label{sec:api-functions}

This section presents the six most important functions in the
\texttt{cardano-license} API, with usage examples.  Every function is
\texttt{async} and must be called within an asyncio event loop.

\subsection{generate\_wallet}

The \texttt{generate\_wallet} function creates a new HD~wallet following the
CIP-1852 derivation standard.  It generates a 24-word BIP-39 mnemonic,
derives payment and stake key pairs using hardened derivation paths
(\texttt{m/1852'/1815'/0'/0/0} for payment, \texttt{m/1852'/1815'/0'/2/0}
for stake), constructs a base address that combines both keys, encrypts the
private keys with AES-256-GCM, saves them to disk, and records the wallet
metadata in the \texttt{blockchain\_wallets} table.

The function accepts a \texttt{wallet\_type} parameter that must be one of
\texttt{authority}, \texttt{licensee}, \texttt{signer}, or
\texttt{observer}.  This type determines the wallet's role in the credential
system and is enforced by a CHECK constraint in the database schema.

\begin{lstlisting}[language=Python,
  caption={Generating an authority wallet.}]
import asyncio
from cardano_license import generate_wallet

async def main():
    wallet = await generate_wallet(
        wallet_type="authority",
        label="state_board_of_engineers",
    )
    print(f"Address:  {wallet['address']}")
    print(f"Mnemonic: {wallet['mnemonic']}")
    # Mnemonic: 24-word BIP-39 phrase (store securely!)
    # Keys saved to ~/.cardano-license/wallets/

asyncio.run(main())
\end{lstlisting}

\subsection{mint\_license\_nft}

The \texttt{mint\_license\_nft} function mints a CIP-25--compliant license
NFT and sends it to a licensee's address.  Under the hood, it loads the
authority wallet's keys from encrypted storage, creates a
\texttt{ScriptPubkey} minting policy restricted to the authority's
verification key hash, builds CIP-25 metadata containing the license
fields, constructs and signs a minting transaction, submits it to the
network via Blockfrost, and records the mint in the
\texttt{blockchain\_licenses} table.

The function validates all required metadata fields before performing any
on-chain operation, failing fast with a \texttt{ValueError} if any field is
missing.  The minting policy ensures that only the authority's signing key
can authorize the creation (or destruction) of tokens under that policy.

\begin{lstlisting}[language=Python,
  caption={Minting a license NFT for a professional engineer.}]
import asyncio
from cardano_license import mint_license_nft

async def main():
    result = await mint_license_nft(
        authority_wallet_label="state_board_of_engineers",
        licensee_address="addr_test1qz...",
        license_metadata={
            "name": "PE License - Jane Smith",
            "license_type": "Professional Engineer",
            "license_number": "PE-2026-00142",
            "jurisdiction": "New York",
            "issued_date": "2026-01-15",
            "expiry_date": "2028-01-15",
            "holder_name": "Jane Smith",
        },
    )
    print(f"Policy ID:  {result['policy_id']}")
    print(f"Token Name: {result['token_name']}")
    print(f"Tx Hash:    {result['tx_hash']}")

asyncio.run(main())
\end{lstlisting}

\subsection{update\_license\_status}

The \texttt{update\_license\_status} function implements CIP-68
authority-controlled revocation.  In the CIP-68 model, the authority holds a
Reference Token (label~100) whose inline datum contains the license status.
To revoke a license, the authority updates this datum to
\texttt{status='revoked'}, requiring only the authority's signature and
\emph{zero interaction} with the licensee's wallet.

This is a critical design property: the licensee's User Token (label~222)
remains in their wallet, but any verifier who checks the Reference Token's
datum will see the revoked status.  The authority retains unilateral control
over credential status without needing to burn or recall tokens from the
licensee.

\begin{lstlisting}[language=Python,
  caption={Revoking a license via CIP-68 datum update.}]
import asyncio
from cardano_license import update_license_status

async def main():
    result = await update_license_status(
        license_id=42,
        new_status="revoked",
        authority_wallet_label="state_board_of_engineers",
    )
    print(f"Old status: {result['old_status']}")
    print(f"New status: {result['new_status']}")
    print(f"Tx Hash:    {result['tx_hash']}")

asyncio.run(main())
\end{lstlisting}

\subsection{sign\_document}

The \texttt{sign\_document} function records a professional's signature on a
document by transferring signature and validity tokens to a contract wallet
address.  The function builds a transaction that transfers one signature
token and one validity token to the specified contract address, attaches the
document's SHA-256 hash as transaction metadata (label~367), and records the
timestamp in the \texttt{blockchain\_signatures} table.

Each signer deposits tokens to an independent UTxO at the contract address,
avoiding the eUTxO concurrency bottleneck that would arise if all signers
tried to spend the same UTxO simultaneously.  This design enables true
parallel signing.

\begin{lstlisting}[language=Python,
  caption={Signing a document with credential tokens.}]
import asyncio
from cardano_license import sign_document

async def main():
    result = await sign_document(
        signer_wallet_label="jane_smith_pe",
        document_hash="a1b2c3d4e5f6...64-char-sha256-hex",
        contract_address="addr_test1qr...",
        license_ref=42,
    )
    print(f"Signature ID: {result['signature_id']}")
    print(f"Tx Hash:      {result['tx_hash']}")
    print(f"Timestamp:    {result['timestamp']}")

asyncio.run(main())
\end{lstlisting}

\subsection{verify\_signature}

The \texttt{verify\_signature} function checks whether a document has been
signed by querying the local signature records and verifying each signer's
validity-token status.  It returns a comprehensive result including the
verification status, a list of all signatures, the signature count, and
work-product information if the document is part of a multi-signer work
product.

The verification process cross-references the
\texttt{blockchain\_signatures} table against the
\texttt{blockchain\_validity\_tokens} table to ensure that each signer held
a valid (unexpired, unrevoked) validity token at the time of signing.

\begin{lstlisting}[language=Python,
  caption={Verifying a document's signatures.}]
import asyncio
from cardano_license import verify_signature

async def main():
    result = await verify_signature(
        contract_address="addr_test1qr...",
        document_hash="a1b2c3d4e5f6...64-char-sha256-hex",
    )
    print(f"Verified:         {result['is_verified']}")
    print(f"Signature count:  {result['signature_count']}")
    for sig in result["signatures"]:
        print(f"  Signer: {sig['signer_address']}")
        print(f"  Time:   {sig['timestamp']}")

asyncio.run(main())
\end{lstlisting}

\subsection{pay\_dues}

The \texttt{pay\_dues} function records a dues payment against a
dues-enforcement contract and renews the associated validity token.  It
validates the payment amount against the contract's
\texttt{annual\_dues\_lovelace} requirement, records the payment in the
\texttt{dues\_payments} table, and updates the validity expiry.

The dues-enforcement system links credential validity to ongoing financial
obligations.  If a licensee fails to pay annual dues within the grace
period, the authority can invoke the \texttt{RevokeValidityRedeemer} to
revoke the validity token on-chain, effectively disabling the credential.

\begin{lstlisting}[language=Python,
  caption={Paying annual dues to maintain credential validity.}]
import asyncio
from cardano_license import pay_dues

async def main():
    result = await pay_dues(
        contract_id=7,
        payer_address="addr_test1qz...",
        payment_lovelace=50_000_000,   # 50 ADA
        new_expiry="2027-03-01T00:00:00Z",
    )
    print(f"Payment ID:  {result['payment_id']}")
    print(f"New Expiry:  {result['new_expiry']}")

asyncio.run(main())
\end{lstlisting}


% ----------------------------------------------------------------------------
\section{Database Schema}
\label{sec:db-schema}

The \texttt{cardano-license} package maintains local state in a SQLite
database with ten tables.  This database is \emph{not} a replacement for the
blockchain---it serves as a local index and cache that enables fast queries
without hitting the Blockfrost API for every operation.  The on-chain state
remains the authoritative source of truth for all token ownership,
transaction finality, and datum values.

Table~\ref{tab:db-schema} summarizes all ten tables.  The
entity-relationship diagram in Figure~\ref{fig:er-diagram} shows the key
foreign-key relationships.

\begin{table}[htbp]
\centering
\caption{Database schema: all ten tables in \texttt{cardano-license}.}
\label{tab:db-schema}
\small
\begin{tabular}{@{}p{3.8cm}p{3.2cm}p{5.5cm}@{}}
\toprule
\textbf{Table} & \textbf{Key Columns} & \textbf{Purpose} \\
\midrule
\texttt{blockchain\_wallets}
  & id, wallet\_type, address, public\_key\_hash, network, label
  & Stores wallet metadata for all four roles (authority, licensee, signer, observer) \\
\addlinespace
\texttt{blockchain\_licenses}
  & id, token\_name, policy\_id, licensee\_address, authority\_address, status
  & Tracks minted license NFTs with metadata, status, and validity dates \\
\addlinespace
\texttt{blockchain\_reference\_tokens}
  & id, license\_id, policy\_id, user\_token\_name, ref\_token\_name, datum\_json, status
  & CIP-68 reference tokens with mutable datum for authority-controlled revocation \\
\addlinespace
\texttt{blockchain\_signature\_tokens}
  & id, policy\_id, token\_name, licensee\_address, quantity, status
  & Signature token inventory per licensee \\
\addlinespace
\texttt{blockchain\_validity\_tokens}
  & id, policy\_id, token\_name, licensee\_address, valid\_until, status
  & Time-bounded validity tokens linked to licenses \\
\addlinespace
\texttt{blockchain\_signatures}
  & id, document\_hash, signer\_address, signature\_tx\_hash, signature\_datum
  & Individual signature records with document hash and timestamp \\
\addlinespace
\texttt{blockchain\_work\_products}
  & id, title, wp\_address, document\_hash, required\_signers\_json, status
  & Multi-signer work products with signature collection tracking \\
\addlinespace
\texttt{blockchain\_minting\_policies}
  & policy\_id, authority\_address, policy\_cbor\_hex, script\_type
  & Minting policy registry (native script and Plutus~V2) \\
\addlinespace
\texttt{dues\_contracts}
  & id, authority\_address, license\_ref, annual\_dues\_lovelace, grace\_period\_slots, status
  & Dues-enforcement contracts linking credential validity to financial obligations \\
\addlinespace
\texttt{dues\_payments}
  & id, contract\_id, payer\_address, amount\_lovelace, new\_expiry, status
  & Payment records against dues contracts with confirmation tracking \\
\bottomrule
\end{tabular}
\end{table}

\begin{figure}[htbp]
\centering
\begin{tikzpicture}[
  entity/.style={guidebox, fill=blue!8, minimum width=3.4cm, minimum height=0.9cm,
    font=\sffamily\scriptsize\bfseries},
  node distance=1.3cm and 2.6cm
]
  % ── Entities ──
  \node[entity] (wallets) {blockchain\_wallets};
  \node[entity, right=2.8cm of wallets] (licenses) {blockchain\_licenses};
  \node[entity, below=1.5cm of licenses] (reftokens) {blockchain\_reference\_tokens};
  \node[entity, below=1.5cm of wallets] (sigtokens) {blockchain\_signature\_tokens};
  \node[entity, below=1.5cm of reftokens] (valtokens) {blockchain\_validity\_tokens};
  \node[entity, below=1.5cm of sigtokens] (signatures) {blockchain\_signatures};
  \node[entity, below=1.5cm of signatures] (workproducts) {blockchain\_work\_products};
  \node[entity, below=1.5cm of valtokens] (policies) {blockchain\_minting\_policies};
  \node[entity, below=1.5cm of workproducts] (contracts) {dues\_contracts};
  \node[entity, below=1.5cm of policies] (payments) {dues\_payments};

  % ── Relationships ──
  \draw[guidearrow] (licenses) -- node[guidelabel, above] {authority\_address} (wallets);
  \draw[guidearrow] (reftokens) -- node[guidelabel, right, text width=1.8cm] {license\_id} (licenses);
  \draw[guidearrow] (sigtokens.east) -- node[guidelabel, above] {license\_ref} (licenses.south west);
  \draw[guidearrow] (valtokens) -- node[guidelabel, right, text width=1.8cm] {license\_ref} (reftokens);
  \draw[guidearrow] (signatures) -- node[guidelabel, above] {signer\_address} (sigtokens);
  \draw[guidearrow] (workproducts) -- node[guidelabel, left, text width=1.5cm] {signatures} (signatures);
  \draw[guidearrow] (contracts.east) -- node[guidelabel, above] {license\_ref} (payments.west);
  \draw[guidearrow] (contracts) -- node[guidelabel, left, text width=1.8cm] {license\_ref} (workproducts);
\end{tikzpicture}
\caption{Entity-relationship diagram for the \texttt{cardano-license} database.
  Arrows indicate referential relationships between tables.}
\label{fig:er-diagram}
\end{figure}

The schema enforces data integrity through CHECK constraints on enumerated
columns.  For example, the \texttt{blockchain\_wallets.wallet\_type} column
is constrained to \texttt{('authority', 'licensee', 'signer', 'observer')},
and the \texttt{blockchain\_licenses.status} column is constrained to
\texttt{('pending', 'active', 'revoked', 'expired', 'burned')}.  These
constraints prevent invalid state transitions at the database level,
providing a safety net independent of application-level validation.


% ----------------------------------------------------------------------------
\section{Configuration}
\label{sec:configuration}

The \texttt{cardano-license} package is configured entirely through
environment variables, following the twelve-factor app methodology.  This
design enables the same codebase to target mainnet, preprod, or preview
networks without code changes.

\begin{table}[htbp]
\centering
\caption{Environment variables for \texttt{cardano-license} configuration.}
\label{tab:env-vars}
\begin{tabular}{@{}p{4.6cm}p{2.8cm}p{5cm}@{}}
\toprule
\textbf{Variable} & \textbf{Default} & \textbf{Purpose} \\
\midrule
\texttt{CARDANO\_NETWORK}
  & \texttt{testnet}
  & Target network: \texttt{mainnet}, \texttt{testnet} (preprod), or \texttt{preview} \\
\addlinespace
\texttt{BLOCKFROST\_PROJECT\_ID}
  & (empty)
  & Blockfrost API key for chain queries and transaction submission \\
\addlinespace
\texttt{CARDANO\_LICENSE\_PASSWORD}
  & (empty)
  & Password for AES-256-GCM wallet key encryption; if empty, prompts interactively \\
\addlinespace
\texttt{CARDANO\_LICENSE\_DIR}
  & \texttt{\textasciitilde/.cardano-license}
  & Base directory for wallets, policies, and database \\
\addlinespace
\texttt{CARDANO\_LICENSE\_DB}
  & \texttt{<DIR>/license.db}
  & SQLite database file path \\
\bottomrule
\end{tabular}
\end{table}

The configuration module (\texttt{config.py}) automatically creates the
required directory structure on import: the base directory, the
\texttt{wallets/} subdirectory for encrypted key files, and the
\texttt{policies/} subdirectory for serialized minting-policy scripts.  This
means that a fresh installation requires no manual setup beyond setting the
environment variables.

For production deployments, the \texttt{CARDANO\_NETWORK} variable should be
set to \texttt{mainnet} and the \texttt{BLOCKFROST\_PROJECT\_ID} should
contain a mainnet API key obtained from the Blockfrost dashboard.  For
development and testing, the default \texttt{testnet} setting targets the
Cardano preprod network, where test ADA is freely available from faucets.


% ----------------------------------------------------------------------------
\section{Testing}
\label{sec:testing}

The \texttt{cardano-license} package includes a comprehensive test suite
with 660~passing tests.  The tests are organized into four test modules
covering the full spectrum of package functionality.

\texttt{test\_license.py} covers wallet management (generation, key
encryption/decryption, label lookup), license NFT operations (minting,
metadata validation, status queries), CIP-68 reference-token management
(store, update, revocation), signature-token operations (mint, balance,
transfer), validity-token operations (mint, check, renew, revoke), document
signing and verification, work-product lifecycle (create, sign, finalize),
and authority registration.

\texttt{test\_contracts.py} covers Plutus~V2 smart-contract operations:
signature-collection validator construction, datum and redeemer
serialization, dues-enforcement contract deployment, payment validation, and
grace-period logic.

\texttt{test\_tx\_utils.py} covers transaction-building utilities: fee
estimation, UTxO selection, mint/transfer/multisig transaction construction,
submission, and confirmation polling.

\texttt{test\_testnet.py} contains integration tests designed to run against
the Cardano preprod testnet with a live Blockfrost connection.  These tests
are marked with \texttt{@pytest.mark.testnet} and are skipped by default in
CI environments where testnet access is unavailable.

The test suite uses \texttt{pytest} with the \texttt{pytest-asyncio} plugin
for async test support.  All database-touching tests use isolated temporary
databases to prevent cross-test contamination.  Running the full suite
requires only:

\begin{lstlisting}[language=bash, caption={Running the test suite.}]
pytest tests/ -v
# 660 passed in ~45s (excluding testnet tests)
\end{lstlisting}


% ============================================================================
%  PART B — THE v4.0 ROADMAP
% ============================================================================

\vspace{1em}
\noindent\rule{\textwidth}{0.8pt}
\begin{center}
  {\Large\sffamily\bfseries Part B --- The v4.0 Roadmap}
\end{center}
\noindent\rule{\textwidth}{0.4pt}
\vspace{0.5em}


% ----------------------------------------------------------------------------
\section{Aiken Smart Contracts}
\label{sec:aiken}

The current implementation uses Plutus~V2 validators written as PyCardano
\texttt{PlutusV2Script} objects with serialized CBOR bytecode.  While
functional, this approach inherits the well-known limitations of the
PlutusTx/Haskell compilation pipeline: large compiled UPLC (Untyped Plutus
Lambda Calculus) output, high CPU and memory execution-unit (exUnit)
consumption, and a steep learning curve for developers unfamiliar with
Haskell.

Aiken is a modern smart-contract language designed specifically for Cardano.
It uses a Rust-like syntax that is far more accessible to the broader
developer community while maintaining the same security guarantees as
Haskell-compiled Plutus scripts.  The Aiken compiler produces significantly
smaller UPLC bytecode---often 5--10$\times$ smaller than equivalent PlutusTx
output---which directly translates to lower exUnit consumption and lower
transaction fees.

The v4.0 roadmap calls for rewriting all three of the credential system's
validators in Aiken: the signature-collection validator (which manages
multi-signer deposits, finalization sweeps, and reclaim operations), the
dues-enforcement validator (which links credential validity to payment
obligations), and the planned validity-consolidation validator (described in
Section~\ref{sec:validity-consolidation}).

The benefits of the Aiken migration are fourfold.  First, the compiled
script sizes will shrink dramatically, reducing the on-chain footprint.
Second, the smaller scripts consume fewer exUnits, lowering per-transaction
costs.  Third, the Rust-like syntax lowers the barrier to community
contributions and third-party auditing.  Fourth, the Aiken toolchain
includes built-in testing and property-based fuzzing, improving validator
correctness assurance.

\begin{figure}[htbp]
\centering
\begin{tikzpicture}[node distance=0.9cm and 1.2cm]
  % ── PlutusTx workflow (left) ──
  \node[guidelabel] at (-3.5, 3.2) {\textbf{Current: PlutusTx Workflow}};

  \node[guidebox, fill=red!6, minimum width=2.8cm] (haskell)
    at (-3.5, 2.2) {Haskell Source\\{\scriptsize\textsf{(.hs files)}}};
  \node[guidebox, fill=red!6, minimum width=2.8cm, below=0.7cm of haskell]
    (ghc) {GHC Compiler\\{\scriptsize\textsf{+ PlutusTx plugin}}};
  \node[guidebox, fill=red!6, minimum width=2.8cm, below=0.7cm of ghc]
    (plutustx) {PlutusTx IR\\{\scriptsize\textsf{(intermediate)}}};
  \node[guidebox, fill=red!6, minimum width=2.8cm, below=0.7cm of plutustx]
    (uplc1) {UPLC Bytecode\\{\scriptsize\textsf{(large output)}}};

  \draw[guidearrow] (haskell) -- (ghc);
  \draw[guidearrow] (ghc) -- (plutustx);
  \draw[guidearrow] (plutustx) -- (uplc1);

  % ── Aiken workflow (right) ──
  \node[guidelabel] at (3.5, 3.2) {\textbf{v4.0: Aiken Workflow}};

  \node[guidebox, fill=green!8, minimum width=2.8cm] (aiken)
    at (3.5, 2.2) {Aiken Source\\{\scriptsize\textsf{(.ak files)}}};
  \node[guidebox, fill=green!8, minimum width=2.8cm, below=0.7cm of aiken]
    (aikencc) {Aiken Compiler\\{\scriptsize\textsf{(Rust-based)}}};
  \node[guidebox, fill=green!8, minimum width=2.8cm, below=0.7cm of aikencc]
    (uplc2) {UPLC Bytecode\\{\scriptsize\textsf{(compact output)}}};

  \draw[guidearrow] (aiken) -- (aikencc);
  \draw[guidearrow] (aikencc) -- (uplc2);

  % ── Comparison annotation ──
  \draw[dashed, thick, gray!50] (0, 3.5) -- (0, -2.2);
  \node[guidelabel, text width=2.5cm, align=center] at (0, -2.8)
    {5--10$\times$ smaller\\output size};
\end{tikzpicture}
\caption{Compilation pipeline comparison: PlutusTx (left) requires three
  stages through the GHC/Haskell toolchain, producing large UPLC.  Aiken
  (right) compiles directly to compact UPLC in a single step.}
\label{fig:aiken-comparison}
\end{figure}


% ----------------------------------------------------------------------------
\section{CIP-33 Reference Scripts}
\label{sec:cip33}

In the current architecture, every transaction that invokes a Plutus
validator must attach the full compiled script bytecode.  For the
signature-collection validator, this adds several kilobytes to each signing
transaction.  For the dues-enforcement validator, it adds similar overhead
to every payment transaction.  This per-transaction overhead increases both
transaction size and fees.

CIP-33 (Reference Scripts) solves this problem by allowing validators to be
deployed to the blockchain \emph{once} as reference scripts.  The deploying
transaction creates a UTxO that carries the compiled script in its output.
Every subsequent transaction that needs to invoke that validator simply
references the UTxO containing the script, rather than attaching the full
bytecode.

The impact on transaction economics is substantial.  A typical
signature-deposit transaction currently weighs approximately 8--12~KB due to
the attached validator script.  With CIP-33 reference scripts, the same
transaction would weigh only 400--800~bytes---the script reference adds just
a 32-byte transaction hash and a UTxO index.  This reduction translates
directly to lower fees: roughly 80--90\% savings on script-bearing
transactions.

The v4.0 implementation plan is straightforward.  First, deploy each
validator (signature collection, dues enforcement, validity consolidation)
as a reference script in a dedicated UTxO.  Second, update the transaction
builders in \texttt{tx\_utils.py} to reference the script UTxO instead of
attaching the full bytecode.  Third, store the reference-script UTxO
identifiers in the \texttt{blockchain\_minting\_policies} table for
retrieval during transaction construction.

\begin{figure}[htbp]
\centering
\begin{tikzpicture}[node distance=1cm and 1.5cm]
  % ── Before: Script Attached ──
  \node[guidelabel] at (-3.8, 2.8) {\textbf{Before: Script Attached}};

  \node[guidebox, fill=red!6, minimum width=3cm, minimum height=1.4cm]
    (txbefore) at (-3.8, 1.2)
    {\textbf{Transaction}\\{\scriptsize\textsf{Inputs + Outputs}}\\
     {\scriptsize\textsf{+ \textbf{Full Script} (8 KB)}}};

  \node[guidebox, fill=gray!10, minimum width=2.6cm, below=0.8cm of txbefore]
    (chain1) {Cardano Ledger\\{\scriptsize\textsf{Validates inline script}}};

  \draw[guidearrow] (txbefore) -- (chain1);

  % ── After: Script Referenced ──
  \node[guidelabel] at (3.8, 2.8) {\textbf{After: Script Referenced}};

  \node[guidebox, fill=green!8, minimum width=3cm, minimum height=1.4cm]
    (txafter) at (3.8, 1.2)
    {\textbf{Transaction}\\{\scriptsize\textsf{Inputs + Outputs}}\\
     {\scriptsize\textsf{+ \textbf{Script Ref} (32 B)}}};

  \node[guidebox, fill=blue!8, minimum width=2.6cm, right=0.6cm of txafter]
    (refutxo) {Reference UTxO\\{\scriptsize\textsf{Contains full}}\\
     {\scriptsize\textsf{script (deployed once)}}};

  \node[guidebox, fill=gray!10, minimum width=2.6cm, below=0.8cm of txafter]
    (chain2) {Cardano Ledger\\{\scriptsize\textsf{Reads script from ref UTxO}}};

  \draw[guidearrow] (txafter) -- (chain2);
  \draw[guidearrow, dashed] (txafter) -- (refutxo);
  \draw[guidearrow] (refutxo) |- (chain2);

  % ── Savings annotation ──
  \node[guidelabel, text width=3cm, align=center] at (0, -1.8)
    {$\sim$90\% reduction\\in transaction size};
\end{tikzpicture}
\caption{CIP-33 Reference Scripts: before (left), each transaction carries
  the full validator bytecode.  After (right), the script is deployed once
  and subsequent transactions reference it by UTxO, dramatically reducing
  size and fees.}
\label{fig:cip33-diagram}
\end{figure}


% ----------------------------------------------------------------------------
\section{Validity Token Consolidation}
\label{sec:validity-consolidation}

The current three-token architecture presents a dual-state consistency risk
that was identified during the v3.0 security audit (see
Section~\ref{sec:audit}).  The problem: the CIP-68 Reference Token's datum
may say ``revoked,'' but the licensee may still hold an unexpired Validity
Token (Token~3) in their wallet.  A na\"ive verifier that checks only the
Validity Token would incorrectly conclude that the credential is still
valid.  Although the system documentation specifies that verifiers must check
\emph{both} the Reference Token datum and the Validity Token, this dual-check
requirement introduces a class of verification errors that a simpler
architecture would eliminate.

The v4.0 solution is to deprecate Token~3 (the standalone Validity Token)
entirely.  Instead, validity information will be encoded directly into the
CIP-68 Reference Token's datum as a \texttt{valid\_until} field containing a
POSIX timestamp.  Any verifier---or any on-chain validator---can read this
datum via a CIP-31 Reference Input (a read-only input that does not consume
the UTxO) and determine both the credential's status and its validity
window from a \emph{single source of truth}.

This consolidation eliminates the dual-state risk entirely.  The authority
controls the Reference Token's datum, and the datum contains both the status
(\texttt{active}/\texttt{revoked}/\texttt{suspended}) and the validity
window (\texttt{valid\_until}).  There is no separate token whose state
could contradict the Reference Token.

The Plutus validator for the new consolidated design is straightforward: it
reads the Reference Token's datum via CIP-31 Reference Input, extracts the
\texttt{valid\_until} field, compares it against the current slot
(converted to POSIX time), and checks the \texttt{status} field---all in a
single on-chain read with no UTxO consumption.

\begin{figure}[htbp]
\centering
\begin{tikzpicture}[node distance=1.2cm and 1.8cm]
  % ── Before: Dual State ──
  \node[guidelabel] at (-4, 3.6) {\textbf{Before: Dual-State Risk}};

  \node[guidebox, fill=blue!8, minimum width=2.8cm] (refbefore)
    at (-4, 2.4) {Reference Token\\{\scriptsize\textsf{datum: status=revoked}}};

  \node[guidebox, fill=orange!8, minimum width=2.8cm, below=0.9cm of refbefore]
    (valtok) {Validity Token\\{\scriptsize\textsf{valid\_until: 2027-03}}};

  \node[guidebox, fill=red!6, minimum width=3cm, below=0.9cm of valtok]
    (conflict) {\textbf{Conflict!}\\{\scriptsize\textsf{Revoked but token says valid}}};

  \draw[guidearrow] (refbefore) -- node[guidelabel, right] {says revoked} (valtok);
  \draw[guidearrow] (valtok) -- node[guidelabel, right] {says valid} (conflict);

  % ── After: Single Source ──
  \node[guidelabel] at (4, 3.6) {\textbf{After: Single Source of Truth}};

  \node[guidebox, fill=green!8, minimum width=3.2cm, minimum height=1.6cm] (refafter)
    at (4, 2.1) {Reference Token\\{\scriptsize\textsf{datum: status=revoked}}\\
     {\scriptsize\textsf{valid\_until: 2027-03}}};

  \node[guidebox, fill=green!8, minimum width=3cm, below=1.2cm of refafter]
    (verifier) {Verifier / Validator\\{\scriptsize\textsf{CIP-31 Reference Input}}\\
     {\scriptsize\textsf{reads datum (no UTxO spent)}}};

  \draw[guidearrow] (refafter) -- node[guidelabel, right] {single read} (verifier);

  % ── Annotation ──
  \draw[dashed, thick, gray!50] (0, 4) -- (0, -1.2);
  \node[guidelabel, text width=2.5cm, align=center] at (0, -1.8)
    {Token~3\\deprecated};
\end{tikzpicture}
\caption{Validity token consolidation: before (left), the Reference Token and
  Validity Token can contradict each other.  After (right), all status and
  validity information resides in the Reference Token's datum, readable via
  CIP-31 Reference Input.}
\label{fig:validity-consolidation}
\end{figure}


% ----------------------------------------------------------------------------
\section{HSM Authority Keys}
\label{sec:hsm}

In the current implementation, authority wallet private keys are encrypted
with AES-256-GCM using a password-derived key (PBKDF2-HMAC-SHA256 with
100\,000 iterations) and stored as files on the server's filesystem.  While
this provides strong cryptographic protection against offline attacks, the
keys are still \emph{software-managed}: they are decrypted into process
memory during signing operations, and a sufficiently privileged attacker
with server access could potentially extract them.

Hardware Security Modules (HSMs) provide a fundamentally stronger security
model.  An HSM is a dedicated cryptographic processor that generates, stores,
and uses keys entirely within a tamper-resistant hardware enclave.  The
critical property is that private keys \emph{never leave the HSM}---not even
during signing.  When the application needs to sign a transaction, it sends
the transaction hash to the HSM, and the HSM returns the signature.  The
private key material is never exposed to the host operating system, the
application process, or any software layer.

For the credential verification system, HSM integration means that the
authority's signing key---the key that controls who can mint license NFTs,
revoke credentials, and deploy smart contracts---would be protected at the
hardware level.  Even if an attacker gained root access to the server, they
could not extract the authority key.  They could potentially \emph{use} the
HSM to sign transactions (if they gained access to the HSM's authentication
credentials), but they could not \emph{steal} the key for use elsewhere.

The v4.0 HSM integration plan targets PKCS\#11-compatible HSMs, which
include cloud-hosted options (AWS CloudHSM, Azure Dedicated HSM, Google
Cloud HSM) and on-premises devices (Thales Luna, Utimaco CryptoServer).
The \texttt{crypto.py} module will be extended with an HSM backend that
implements the same \texttt{encrypt\_key}/\texttt{decrypt\_key} interface,
allowing the rest of the codebase to use HSM-backed keys without
modification.


% ----------------------------------------------------------------------------
\section{KERI / Multi-Sig Upgrade}
\label{sec:keri-multisig}

The current system architecture concentrates all authority operations behind
a single signing key.  Whoever controls that key can mint credentials,
revoke licenses, deploy smart contracts, and collect dues payments.  This
single-key design was identified as finding AS-6 in the v3.0 security audit:
a single point of failure that, if compromised, could undermine the entire
credential system for that authority.

The v4.0 roadmap addresses this with two alternative approaches to
distributed governance, either of which eliminates the single-key risk.

\textbf{Option A: KERI (Key Event Receipt Infrastructure).}
KERI, standardized through the Veridian project, provides a decentralized
key-management framework where key authority is established through a
cryptographically linked sequence of key events (inception, rotation,
delegation) rather than through a single static key.  In a KERI-based
authority model, the credential-issuing authority would be identified by its
KERI Autonomic Identifier (AID), and key rotation events would allow the
authority to update its signing keys without invalidating previously issued
credentials.  This provides resilience against key compromise: if a key is
suspected of being compromised, the authority can rotate to a new key pair,
and all future operations use the new key while past operations remain
verifiable through the key-event log.

\textbf{Option B: Cardano ScriptNofK Multi-Signature.}
Cardano's native scripting system supports $N$-of-$K$ multi-signature
policies through the \texttt{ScriptNofK} constructor.  A 3-of-5
configuration, for example, would require any three out of five designated
board members to co-sign every authority operation.  This distributes trust
across multiple individuals and eliminates the single-key bottleneck.

In practical terms, a 3-of-5 multi-sig policy means that minting a new
license NFT requires three board members to independently sign the
transaction.  Revoking a credential requires three signatures.  Deploying a
new smart contract requires three signatures.  No single individual---and no
single compromised key---can perform any authority action alone.

The choice between KERI and multi-sig depends on the deploying
organization's requirements.  KERI provides more sophisticated key
management (rotation, delegation, recovery) but requires infrastructure for
maintaining the key-event log.  Multi-sig is simpler to implement using
Cardano's native scripting but does not support key rotation without
creating a new policy (and thus new policy ID).  Both approaches achieve the
primary goal: eliminating the single-key centralization risk identified in
AS-6.


% ----------------------------------------------------------------------------
\section{Security Audit Results}
\label{sec:audit}

The credential verification system underwent three rounds of architectural
security audit in February 2026, progressing from initial finding
identification through remediation verification to final clearance.

\subsection{v3.0 Audit: Initial Findings}

The first audit round examined the v3.0 architecture and identified five
findings across three severity levels:

\begin{itemize}[nosep]
  \item \textbf{2 Critical findings:}
    \begin{enumerate}[nosep]
      \item \emph{Single-key authority risk (AS-6):} All authority operations
        depend on a single signing key.  Compromise of this key would allow
        an attacker to mint fraudulent credentials, revoke legitimate ones,
        or deploy malicious smart contracts.
      \item \emph{Dual-state validity inconsistency:} The CIP-68 Reference
        Token datum and the standalone Validity Token can hold contradictory
        state, creating a class of verification errors.
    \end{enumerate}
  \item \textbf{1 High finding:}
    \begin{enumerate}[nosep]
      \item \emph{Software-managed authority keys:} Private keys decrypted
        into process memory during signing operations are vulnerable to
        memory-extraction attacks by privileged processes.
    \end{enumerate}
  \item \textbf{2 Medium findings:}
    \begin{enumerate}[nosep]
      \item \emph{Per-transaction script attachment overhead:} Attaching full
        validator bytecode to every transaction increases fees and exposes the
        compiled script to analysis.
      \item \emph{PlutusTx compilation inefficiency:} The Haskell-based
        compilation pipeline produces unnecessarily large UPLC output.
    \end{enumerate}
\end{itemize}

\subsection{v3.1 Audit: Remediation Verification}

The second audit round verified that all five findings from the v3.0 audit
had been addressed with concrete remediation plans:

\begin{itemize}[nosep]
  \item AS-6 (single-key risk): addressed by the KERI/multi-sig upgrade
    (Section~\ref{sec:keri-multisig})
  \item Dual-state inconsistency: addressed by validity-token consolidation
    (Section~\ref{sec:validity-consolidation})
  \item Software-managed keys: addressed by HSM integration
    (Section~\ref{sec:hsm})
  \item Script attachment overhead: addressed by CIP-33 reference scripts
    (Section~\ref{sec:cip33})
  \item PlutusTx inefficiency: addressed by Aiken migration
    (Section~\ref{sec:aiken})
\end{itemize}

\subsection{Final Audit: Testnet Clearance}

The third and final audit round confirmed that the remediation plans were
architecturally sound and that the v3.1 implementation (with the v4.0
roadmap items scheduled for implementation) was cleared for testnet
deployment.  All three audit reports are published in the project
repository.

The audit progression demonstrates a disciplined approach to security
engineering: identify weaknesses, design mitigations, verify the mitigations
address the root causes, and document everything.  The v4.0 roadmap items
described in this chapter are the direct result of this audit process.


% ============================================================================
%  PART C — GLOSSARY
% ============================================================================

\vspace{1em}
\noindent\rule{\textwidth}{0.8pt}
\begin{center}
  {\Large\sffamily\bfseries Part C --- Glossary}
\end{center}
\noindent\rule{\textwidth}{0.4pt}
\vspace{0.5em}

\section{Glossary of Terms}
\label{sec:glossary}

The following glossary defines every significant technical term used
throughout this study guide, organized alphabetically.

\begin{description}[style=nextline, leftmargin=2.5cm, labelwidth=2.3cm]

\item[ADA]
  The native cryptocurrency of the Cardano blockchain.  One ADA equals
  1\,000\,000 lovelace.  ADA is used to pay transaction fees, satisfy
  minimum UTxO requirements, and transfer value between addresses.

\item[AES-256-GCM]
  Advanced Encryption Standard with 256-bit keys in Galois/Counter Mode.
  An authenticated encryption algorithm that provides both confidentiality
  and integrity.  Used by \texttt{cardano-license} to encrypt wallet private
  keys at rest, with key derivation via PBKDF2-HMAC-SHA256 (100\,000
  iterations).

\item[Aiken]
  A modern smart-contract language for Cardano with Rust-like syntax.
  Compiles directly to UPLC bytecode, producing significantly smaller output
  than the PlutusTx/Haskell pipeline.  Planned for v4.0 to replace
  Haskell-compiled validators.

\item[Asset Name]
  The human-readable identifier for a native token within a minting policy.
  Combined with the Policy~ID to form a globally unique token identifier.
  In the credential system, asset names encode the license holder's
  identifier (e.g., \texttt{PE-2026-00142}).

\item[Base Address]
  A Cardano address that combines both a payment credential (derived from
  the payment key) and a staking credential (derived from the stake key).
  This is the standard address format used by the credential system, enabling
  both spending and staking from the same address.

\item[BIP-39]
  Bitcoin Improvement Proposal~39: a standard for generating deterministic
  wallet seeds from mnemonic word sequences.  The credential system uses
  24-word BIP-39 mnemonics as the root entropy for all wallet derivation.

\item[BIP-44]
  Bitcoin Improvement Proposal~44: a standard for hierarchical deterministic
  (HD) wallet derivation paths.  Superseded by CIP-1852 for Cardano-specific
  derivation, but the underlying multi-level derivation concept remains the
  same.

\item[Blake2b]
  A cryptographic hash function used extensively in Cardano for address
  derivation, script hashing, and transaction identification.  Cardano uses
  Blake2b-224 (28~bytes) for key hashes and Blake2b-256 (32~bytes) for
  transaction hashes.

\item[Blockfrost]
  A hosted API service that provides access to the Cardano blockchain without
  running a full node.  Supports UTxO queries, transaction submission,
  metadata retrieval, and block monitoring.  Used by \texttt{cardano-license}
  as the default chain backend.

\item[CIP-25]
  Cardano Improvement Proposal~25: the metadata standard for NFTs on
  Cardano.  Defines how token metadata (name, image, description, attributes)
  is attached to minting transactions under metadata label~721.  Used for
  license NFT metadata in the credential system.

\item[CIP-30]
  Cardano Improvement Proposal~30: the wallet connector standard that defines
  a JavaScript API for web applications to interact with Cardano wallets
  (e.g., Nami, Eternl).  Enables browser-based credential verification
  interfaces.

\item[CIP-31]
  Cardano Improvement Proposal~31: Reference Inputs.  Allows transactions to
  read the datum of a UTxO without consuming it.  Critical for the
  validity-token consolidation design, where verifiers read the Reference
  Token's datum without spending the authority's UTxO.

\item[CIP-32]
  Cardano Improvement Proposal~32: Inline Datums.  Allows datums to be
  stored directly in transaction outputs rather than as hashes.  Eliminates
  the need for datum witnesses in spending transactions and enables CIP-31
  reference-input reads.

\item[CIP-33]
  Cardano Improvement Proposal~33: Reference Scripts.  Allows Plutus
  scripts to be deployed once on-chain and referenced by subsequent
  transactions, eliminating the need to attach full script bytecode to every
  transaction.  See Section~\ref{sec:cip33}.

\item[CIP-68]
  Cardano Improvement Proposal~68: an NFT standard that splits each token
  into a Reference Token (label~100, held by the issuer, mutable datum) and a
  User Token (label~222, held by the recipient, immutable).  The credential
  system uses CIP-68 for authority-controlled revocation: the authority
  updates the Reference Token's datum without touching the licensee's wallet.

\item[CIP-1852]
  Cardano Improvement Proposal~1852: the HD wallet derivation standard for
  Cardano.  Defines derivation paths using purpose~1852, coin type~1815
  (Cardano), and separates payment keys (role~0) from stake keys (role~2).
  Path format: \texttt{m/1852'/1815'/0'/\{role\}/\{index\}}.

\item[Crypto-shredding]
  A data-destruction technique where encrypted data is rendered permanently
  inaccessible by destroying the encryption key rather than the ciphertext.
  Used in the credential system's GDPR compliance design: personal data
  stored off-chain under AES-256-GCM can be ``erased'' by deleting the
  encryption key, while the on-chain attestation hash remains.

\item[Datum]
  Data attached to a UTxO on the Cardano blockchain.  In the eUTxO model,
  datums serve as the state of smart contracts.  The credential system uses
  datums to store license status, validity windows, signer deposits, and
  dues-contract parameters.  See also: Inline Datum, CIP-32.

\item[DID]
  Decentralized Identifier: a W3C standard for self-sovereign digital
  identities.  DIDs are URI-formatted identifiers that can be resolved to
  DID Documents containing public keys, service endpoints, and
  authentication methods.  Related to KERI's approach to decentralized
  identity.

\item[eIDAS]
  Electronic Identification, Authentication and Trust Services: the European
  Union regulation establishing a legal framework for electronic signatures,
  seals, timestamps, and identification.  eIDAS~2.0 introduces the European
  Digital Identity Wallet.  The credential system's on-chain signatures align
  with eIDAS requirements for advanced electronic signatures.

\item[ESIGN Act]
  The Electronic Signatures in Global and National Commerce Act (2000): a
  United States federal law that grants electronic signatures the same legal
  standing as handwritten signatures.  Together with UETA, establishes the
  legal foundation for blockchain-based professional credential signatures.

\item[eUTxO]
  Extended Unspent Transaction Output: Cardano's transaction model that
  extends Bitcoin's UTxO model with datums (state), redeemers (actions), and
  script contexts (validators).  The eUTxO model provides deterministic
  transaction validation and enables the credential system's per-signer
  deposit design for concurrent signing.

\item[Epoch]
  A time period in the Cardano blockchain, lasting approximately 5~days
  (432\,000 slots on mainnet).  Epochs are relevant to the credential system
  for stake-reward distribution and protocol-parameter updates that affect
  transaction fees and minUTxO values.

\item[exUnits]
  Execution units: a two-dimensional measure of Plutus script resource
  consumption on Cardano, comprising CPU steps and memory units.  Each
  transaction has maximum exUnit budgets.  Smaller exUnit consumption (e.g.,
  from Aiken-compiled scripts) translates to lower fees and higher
  throughput.

\item[GDPR]
  General Data Protection Regulation: the European Union's data-privacy law
  that includes the ``right to erasure'' (Article~17).  The credential
  system's off-chain personal-data storage with crypto-shredding is designed
  to satisfy GDPR requirements while maintaining on-chain attestation
  immutability.

\item[Halo2]
  A zero-knowledge proof system based on polynomial commitments, developed by
  the Electric Coin Company.  Relevant to future privacy-preserving
  credential verification where a holder can prove possession of a valid
  credential without revealing the credential's details.

\item[HD Wallet]
  Hierarchical Deterministic wallet: a wallet architecture where a single
  seed (mnemonic) can derive an unlimited tree of key pairs using
  deterministic derivation paths (BIP-32/BIP-44/CIP-1852).  The credential
  system uses HD wallets for all participant roles.

\item[HSM]
  Hardware Security Module: a dedicated cryptographic processor that
  generates, stores, and uses keys within a tamper-resistant hardware
  enclave.  Private keys never leave the HSM.  Planned for v4.0 to protect
  authority signing keys.  See Section~\ref{sec:hsm}.

\item[Inline Datum]
  A datum stored directly in a transaction output (enabled by CIP-32),
  rather than as a hash that must be provided as a witness during spending.
  Inline datums enable CIP-31 reference-input reads, which are essential for
  the validity-token consolidation design.

\item[KERI]
  Key Event Receipt Infrastructure: a decentralized key-management protocol
  standardized through the Veridian project.  KERI establishes key authority
  through a cryptographically linked sequence of key events (inception,
  rotation, delegation, interaction).  Planned for v4.0 as one option for
  distributed authority governance.  See Section~\ref{sec:keri-multisig}.

\item[Lovelace]
  The smallest unit of ADA.  One ADA equals 1\,000\,000 lovelace.  Named
  after Ada Lovelace.  All on-chain amounts in the credential system are
  denominated in lovelace.

\item[minUTxO]
  The protocol-enforced minimum amount of ADA that must accompany any UTxO.
  The value depends on the UTxO's size (datum size, number of native tokens).
  The credential system uses a conservative default of 2~ADA
  (2\,000\,000~lovelace) for token-bearing outputs.

\item[Minting Policy]
  A script that governs the creation and destruction of native tokens on
  Cardano.  The credential system uses authority-key-restricted
  \texttt{ScriptPubkey} policies (Native Script) that permit only the
  authority's signing key to mint or burn tokens under that policy.

\item[Native Script]
  A simple scripting language on Cardano (predating Plutus) that supports
  key-signature requirements, time locks, and Boolean combinations
  (all-of, any-of, $N$-of-$K$).  The credential system uses Native Scripts
  for minting policies.  Multi-sig governance would use
  \texttt{ScriptNofK}.

\item[NFT]
  Non-Fungible Token: a unique token with quantity~1.  On Cardano, NFTs are
  native multi-asset tokens with an enforced quantity of one.  License NFTs
  in the credential system are NFTs carrying CIP-25 metadata that attests to
  a professional's credential.

\item[Ouroboros]
  The family of proof-of-stake consensus protocols used by Cardano.
  Ouroboros Praos (the current production variant) provides probabilistic
  finality and is formally verified for security properties.  Relevant to
  the credential system because consensus determines when minted credentials
  become final.

\item[PKH]
  Public Key Hash: the Blake2b-224 hash of a verification (public) key.  Used
  as the payment credential in Cardano addresses and as the identity
  parameter in minting policies.  The credential system's authority-restricted
  policies require the authority's PKH to authorize minting.

\item[Plutus]
  The smart-contract platform on Cardano, based on the eUTxO model.  Plutus
  scripts (validators) run on-chain to enforce spending conditions.  The
  credential system uses Plutus~V2 for the signature-collection validator and
  the dues-enforcement contract.

\item[Policy ID]
  The hash of a minting policy script.  Serves as the unique, immutable
  identifier for all tokens minted under that policy.  In the credential
  system, the Policy~ID cryptographically binds every license NFT,
  signature token, and validity token to the authority that created the
  minting policy.

\item[PyCardano]
  A Python library for building Cardano transactions, managing wallets, and
  interacting with the blockchain.  Provides HD~wallet generation, Plutus~V2
  script support, and CIP-1852 derivation.  The primary Cardano SDK used by
  \texttt{cardano-license}.

\item[Redeemer]
  Data provided by a transaction to a Plutus validator, specifying the
  intended action.  In the credential system, redeemers include
  \texttt{CollectRedeemer} (deposit a signature), \texttt{FinalizeRedeemer}
  (sweep all deposits), \texttt{ReclaimRedeemer} (reclaim a deposit),
  \texttt{PayDuesRedeemer} (pay dues), and \texttt{RevokeValidityRedeemer}
  (revoke for non-payment).

\item[Reference Input]
  A transaction input (enabled by CIP-31) that reads a UTxO's datum without
  consuming the UTxO.  Essential for the validity-consolidation design,
  where verifiers read the Reference Token's datum without spending the
  authority's UTxO.  See CIP-31.

\item[Reference Script]
  A Plutus script stored on-chain in a UTxO (enabled by CIP-33) that can be
  referenced by subsequent transactions instead of being attached inline.
  Reduces transaction size and fees.  See Section~\ref{sec:cip33}.

\item[Reference Token]
  In the CIP-68 standard, the token with label~100 that is held by the
  issuing authority and carries a mutable inline datum.  The credential
  system uses Reference Tokens for authority-controlled status management:
  the authority updates the datum to revoke or suspend a credential without
  interacting with the licensee's wallet.

\item[Script Context]
  The information available to a Plutus validator during execution: the
  transaction being validated (inputs, outputs, minting, fees, validity
  range, signatories), the script's own purpose (spending, minting,
  certifying), and the datum and redeemer.  Validators use the Script Context
  to enforce spending conditions.

\item[Slot]
  The fundamental time unit on the Cardano blockchain, lasting 1~second on
  mainnet.  Slots are used in validity ranges (\texttt{InvalidBefore},
  \texttt{InvalidHereAfter}) and in time-locked minting policies.  The
  credential system converts between POSIX timestamps and slot numbers for
  validity-window enforcement.

\item[Stake Key]
  The key pair used for staking operations in a Cardano HD~wallet.  Derived
  from path \texttt{m/1852'/1815'/0'/2/0}.  Combined with the payment key to
  form a base address.  The credential system derives stake keys for all
  wallet types but uses them primarily for address construction.

\item[UETA]
  Uniform Electronic Transactions Act (1999): a U.S. model law (adopted by
  47 states) that establishes the legal validity of electronic records and
  signatures.  Together with the ESIGN Act, provides the legal foundation
  for blockchain-based credential signatures.

\item[UPLC]
  Untyped Plutus Lambda Calculus: the low-level execution language for
  Plutus smart contracts on Cardano.  Both PlutusTx (Haskell) and Aiken
  compile down to UPLC bytecode.  Smaller UPLC output means lower exUnit
  consumption and lower fees.

\item[User Token]
  In the CIP-68 standard, the token with label~222 that is held by the
  recipient (licensee).  The User Token is immutable and serves as the
  holder's proof of credential issuance.  Its counterpart, the Reference
  Token (label~100), is held by the authority and carries the mutable status
  datum.

\item[UTxO]
  Unspent Transaction Output: the fundamental unit of value on the Cardano
  blockchain.  Each UTxO contains an address, a value (ADA and/or native
  tokens), and optionally a datum.  Spending a UTxO consumes it entirely and
  creates new UTxOs.  The credential system's per-signer deposit design
  creates independent UTxOs to avoid concurrency conflicts.

\item[Validator]
  A Plutus script that governs the conditions under which a UTxO can be
  spent.  The credential system uses two validators: the
  signature-collection validator (manages multi-signer deposits and
  finalization) and the dues-enforcement validator (links credential
  validity to payment obligations).

\item[Veridian]
  The open-source project (formerly under the Trust over IP Foundation) that
  develops the KERI protocol and ACDC (Authentic Chained Data Container)
  standards for decentralized identity and verifiable credentials.  The v4.0
  roadmap considers Veridian/KERI as one option for distributed authority
  governance.

\end{description}
